% =========================================================
% IMPLEMENTACIÓN DEL AGENTE CONVERSACIONAL (tarea O3-05)
% QUÉ VA: el agente híbrido en funcionamiento (servicio Python + interfaz
%         interna de la API; el modelo de lenguaje solo interpreta), siguiendo los cuatro grupos del flujo
%         to-be (proceso-tobe.tex): acceso y consulta general, solicitud de
%         cita, confirmación (recordatorio) y derivación.
%         Evidencia principal = capturas de conversaciones reales de WhatsApp
%         con datos ficticios (número del desarrollador, entorno de desarrollo).
% CONTENIDO:
%   1. Configuración en Meta: app, número de prueba (desarrollo) y número
%      dedicado (producción), webhook verificado por HTTPS, sin verificación
%      de la empresa (coherente con despliegue.tex).
%   1b. Proveedor y modelo de lenguaje elegidos (criterios del Obj. 2: interfaz
%      común, sin entrenamiento con los datos, costo por mensaje) y costo real.
%      Evaluación formal del intérprete (precisión por intención): recomendación.
%   2. Plantillas aprobadas: tabla corta (nombre funcional | uso | botones).
%      El recordatorio usa Confirmar / Cancelar / Reprogramar.
%   3. Capturas: (a) aviso de privacidad y menú, (b) agendar, (c) recordatorio
%      con botones, (d) derivación. Pueden ir en una sola figura 2x2.
%      Incluir al menos un mensaje en texto libre interpretado por el modelo.
%   4. Efecto en la agenda: la cita creada por el agente aparece en la agenda
%      de la secretaria con origen "agente" (enlaza con modulos-sistema.tex).
%   5. Correspondencia con el diagrama de estados (fig:ob2_estados_conversacion):
%      1-2 líneas, sin repetir el diagrama.
% NO VA: clases, código del webhook, tokens o identificadores de Meta (secretos).
% EXTENSIÓN: 1 tabla + 1 figura compuesta + 2 párrafos cortos.
% =========================================================

\subsubsection{Implementación del agente conversacional}

El agente se implementó según la arquitectura híbrida del diseño (Figuras~\ref{fig:ob2_estilo_agente} y \ref{fig:ob2_c4_componentes}): un servicio propio en Python con la máquina de estados, el intérprete de mensajes y los recordatorios, y una interfaz interna en la API que aplica las reglas de la agenda y es la única que escribe. Ambos están desplegados en el servidor; la conexión con WhatsApp Cloud API queda pendiente de las credenciales de Meta. \textit{[Actualizar al conectar: número de prueba en desarrollo, número dedicado en producción, proveedor y modelo de lenguaje elegidos.]}

La Tabla~\ref{tab:ob3_agente} resume lo construido por componente, y la Tabla~\ref{tab:ob3_plantillas}, las plantillas que requieren aprobación de Meta.

\begingroup
\fontsize{9}{11}\selectfont
\renewcommand{\arraystretch}{1.2}
\setlength{\LTpre}{0pt}\setlength{\LTpost}{6pt}
\setlength{\tabcolsep}{4pt}
\begin{xltabular}{\linewidth}{|>{\raggedright\arraybackslash}p{3.4cm}|>{\raggedright\arraybackslash}X|>{\raggedright\arraybackslash}p{2.2cm}|}
\caption{Componentes del agente implementados}
\label{tab:ob3_agente} \\
\hline
\textbf{Componente} & \textbf{Implementación} & \textbf{Requerimiento} \\ \hline
\endfirsthead
\hline
\textbf{Componente} & \textbf{Implementación} & \textbf{Requerimiento} \\ \hline
\endhead
Aviso y menú & Aviso con Acepto y No acepto; menú como lista interactiva & RF09, RF10 \\ \hline
Intérprete de mensajes & Modelo de lenguaje por API con respuesta estructurada y validada; sin respuesta válida, menú & RF10, RF13 \\ \hline
Agendar, reprogramar y cancelar & Especialidad, profesional, día y hora; escritura solo tras \textit{Confirmar}; límites aplicados por la API & RF11, RF12 \\ \hline
Urgencias & Primeros horarios de hoy o mañana y aviso de emergencia, sin consejo médico & RF11 \\ \hline
Recordatorios & Tarea programada que envía la plantilla del día anterior y registra la respuesta & RF14 \\ \hline
Derivación & Enlace al WhatsApp de recepción; el agente deja de responder a ese número por unas horas & RF13 \\ \hline
\end{xltabular}
\endgroup

\begingroup
\fontsize{9}{11}\selectfont
\renewcommand{\arraystretch}{1.2}
\setlength{\LTpre}{0pt}\setlength{\LTpost}{6pt}
\setlength{\tabcolsep}{4pt}
\begin{xltabular}{\linewidth}{|>{\raggedright\arraybackslash}p{3.2cm}|>{\raggedright\arraybackslash}X|>{\raggedright\arraybackslash}p{4.0cm}|}
\caption{Plantillas de mensaje sujetas a aprobación de Meta}
\label{tab:ob3_plantillas} \\
\hline
\textbf{Plantilla} & \textbf{Uso} & \textbf{Opciones} \\ \hline
\endfirsthead
\hline
\textbf{Plantilla} & \textbf{Uso} & \textbf{Opciones} \\ \hline
\endhead
Recordatorio de cita & Día previo a la cita (RF14) & Confirmar, Cancelar, Reprogramar \\ \hline
\end{xltabular}
\endgroup

La Figura~\ref{fig:ob3_conversaciones} muestra los cuatro grupos del flujo propuesto (Figuras~\ref{fig:ob2_tobe_comunicacion_p1} y \ref{fig:ob2_tobe_comunicacion_p2}) en conversaciones de prueba.

\begin{figure}[H]
\centering
\IfFileExists{img/Ob3_Agente_Conversaciones.png}{\includegraphics[width=\linewidth]{img/Ob3_Agente_Conversaciones.png}}{\fbox{\parbox{0.9\linewidth}{\centering Pendiente: img/Ob3\_Agente\_Conversaciones.png (2x2: aviso y menú, agendar con texto libre, recordatorio, derivación)}}}
\caption{Conversaciones con el agente: (a) aviso de privacidad y menú, (b) agendamiento, (c) recordatorio, (d) derivación}
\label{fig:ob3_conversaciones}
\end{figure}

\textit{[2--3 líneas tras la conexión: la cita agendada por el agente aparece en la agenda de la secretaria con su origen; las consultas clínicas o de saldo se derivan a recepción (RF13); el agente no devuelve datos clínicos (RNF03).]}
